% ============================================================
% FORM-AC-01 - Analisis de caso: fallo real en un sistema distribuido
% Caso: GitHub, 21-22 oct. 2018 (clusteres MySQL)
% Compilar: pdflatex main.tex -> biber main -> pdflatex main.tex (x2)
% BORRADOR INICIAL: los textos en esta plantilla son un punto de partida con
% hechos ya verificados contra el informe oficial. Cada integrante debe
% reescribir su apartado con sus propias palabras y completar lo que esta
% entre [corchetes]. Declarar el uso de IA al final (honestamente).
% ============================================================
\documentclass[11pt,a4paper]{article}

\usepackage[utf8]{inputenc}
\usepackage[T1]{fontenc}
\usepackage[spanish,es-tabla,es-noquoting]{babel}
\usepackage{mathptmx}                 % Times
\usepackage[margin=2.54cm]{geometry}
\usepackage{setspace}\singlespacing
\usepackage{microtype}
\usepackage{booktabs,tabularx,array}
\usepackage{graphicx}
\usepackage{xcolor}
\usepackage{tikz}
\usetikzlibrary{babel}
\usepackage{enumitem}
\usepackage{fancyhdr}
\usepackage{titlesec}
\usepackage[hidelinks]{hyperref}
\usepackage[backend=biber,style=ieee,sorting=none]{biblatex}
\addbibresource{referencias.bib}

% Formato: justificado, numero de pagina abajo a la derecha
\pagestyle{fancy}\fancyhf{}\renewcommand{\headrulewidth}{0pt}
\fancyfoot[R]{\thepage}
\titleformat{\section}{\normalfont\bfseries\fontsize{12}{14}\selectfont}{\thesection.}{0.5em}{}
\titlespacing*{\section}{0pt}{8pt}{3pt}
\setlist{nosep,leftmargin=1.2em}
\setlength{\parskip}{2pt}
\renewcommand{\arraystretch}{1.08}

\newcommand{\todo}[1]{\textcolor{red}{[#1]}}

\begin{document}

% ------------------------- CARATULA -------------------------
\begin{titlepage}
\centering
{\large\bfseries UNIVERSIDAD T\'ECNICA ESTATAL DE QUEVEDO}\\[4pt]
Facultad de Ciencias de la Computaci\'on\\
Carrera de Software\\[2cm]
{\Large\bfseries An\'alisis de caso: fallo real en un sistema distribuido}\\[6pt]
{\large FORM-AC-01 --- Incidente de GitHub, 21--22 de octubre de 2018}\\[2cm]
\begin{tabular}{@{}ll@{}}
Asignatura: & Aplicaciones Distribuidas (ISR-701)\\
Docente: & Gleiston C. Guerrero-Ulloa, PhD\\
Caso asignado: & GitHub (2018), clusteres MySQL\\
Equipo: & \todo{nombre del equipo}\\
Integrantes: & \todo{Integrante 1}\\ & \todo{Integrante 2}\\ & \todo{Integrante 3}\\ & \todo{Integrante 4}\\
Paralelo: & \todo{paralelo}\\
Fecha: & \todo{fecha de la versi\'on}\\
Repositorio: & \small\texttt{\todo{URL completa del repositorio}}\\
\end{tabular}
\end{titlepage}

% ------------------------- CONTENIDO (max. 2 paginas) -------------------------
\section{Descripci\'on del incidente}
El 21 y 22 de octubre de 2018, GitHub experiment\'o una degradaci\'on del servicio durante 24~h 11~min que afect\'o principalmente a los cl\'usteres MySQL de metadatos de la plataforma (\textit{issues}, \textit{pull requests} y autenticaci\'on) \cite{github2018}. El incidente se origin\'o por una p\'erdida de conectividad de aproximadamente 43~s durante un mantenimiento de red, lo que desencaden\'o un cambio de topolog\'ia y escrituras no replicadas entre regiones. Aunque no hubo p\'erdida de datos de usuarios, se limitaron temporalmente servicios como \textit{webhooks} y GitHub Pages mientras se priorizaba la integridad de la informaci\'on \cite{github2018}. A continuaci\'on, se presenta la l\'inea de tiempo con los principales acontecimientos (horas en UTC).

\begingroup\small
\renewcommand{\arraystretch}{1.2}
\begin{tabularx}{\linewidth}{@{}l l X@{}}
\toprule
\textbf{Fecha} & \textbf{UTC} & \textbf{Evento} \cite{github2018}\\
\midrule
\multicolumn{3}{@{}l}{\textit{Fase 1. Detecci\'on y escalamiento}}\\
\addlinespace[2pt]
21 oct & 22:52 & \textbf{P\'erdida de conectividad ($\sim$43~s) por mantenimiento de red;} Orchestrator traslada el liderazgo y las escrituras a la costa oeste. Al restaurarse la conexi\'on, existen escrituras divergentes entre regiones.\\
21 oct & 22:54 & Los sistemas de monitoreo generan alertas e inicia la respuesta del equipo de ingenier\'ia.\\
21 oct & 23:02 & Se identifica una topolog\'ia inesperada de bases de datos, con primarios \'unicamente en la costa oeste.\\
21 oct & 23:07 & Se bloquean manualmente las herramientas de despliegue para evitar cambios adicionales.\\
21 oct & 23:09--23:13 & El incidente se escala: el estado pasa a amarillo (23:09) y luego a \textbf{rojo} (23:11--23:13).\\
\addlinespace
\multicolumn{3}{@{}l}{\textit{Fase 2. Contenci\'on y plan de recuperaci\'on}}\\
\addlinespace[2pt]
21 oct & 23:13 & Se confirma el impacto en m\'ultiples cl\'usteres; ante escrituras acumuladas en la costa oeste y datos no replicados, se prioriza la integridad de la informaci\'on sobre un retorno inmediato.\\
21 oct & 23:19 & Se pausan \textit{webhooks} y compilaciones de GitHub Pages para reducir escrituras y proteger la consistencia de los datos.\\
22 oct & 00:05 & Se define el plan de recuperaci\'on: restaurar desde copias de seguridad, sincronizar r\'eplicas y restablecer una topolog\'ia estable.\\
\addlinespace
\multicolumn{3}{@{}l}{\textit{Fase 3. Restauraci\'on y sincronizaci\'on}}\\
\addlinespace[2pt]
22 oct & 00:41 & Inicia la restauraci\'on de los cl\'usteres afectados desde respaldos remotos, mientras se buscan formas de acelerar el proceso.\\
22 oct & 06:51 & Varios cl\'usteres se restauran en la costa este y comienzan a sincronizar datos desde la costa oeste; algunas operaciones presentan latencia.\\
22 oct & 07:46 & GitHub publica una actualizaci\'on en su blog; la comunicaci\'on se retrasa por la suspensi\'on de GitHub Pages.\\
22 oct & 11:12 & \textbf{Los primarios vuelven a la costa este,} mejorando el rendimiento, aunque persisten r\'eplicas con horas de retraso e inconsistencias visibles.\\
22 oct & 13:15 & Ante el aumento de tr\'afico y el retraso de las r\'eplicas, se aprovisionan r\'eplicas de lectura adicionales para distribuir la carga y acelerar la sincronizaci\'on.\\
22 oct & 16:24 & Las r\'eplicas se sincronizan y se restablece la topolog\'ia original; el servicio permanece en estado rojo mientras se procesa el trabajo acumulado.\\
\addlinespace
\multicolumn{3}{@{}l}{\textit{Fase 4. Procesamiento de la cola y cierre}}\\
\addlinespace[2pt]
22 oct & 16:45 & Se reportan m\'as de 5 millones de eventos de \textit{webhooks} y 80\,000 compilaciones de Pages en cola; cerca de 200\,000 \textit{webhooks} expiran por TTL y se ajusta la configuraci\'on.\\
22 oct & 23:03 & \textbf{Se completa el procesamiento y la verificaci\'on de los sistemas; el estado del sitio vuelve a verde y finaliza el incidente.}\\
\bottomrule
\end{tabularx}
\endgroup

\newpage

\section{Causas ra\'iz}
\begin{figure}[h]
\centering
\resizebox{\linewidth}{!}{%
\begin{tikzpicture}[font=\scriptsize, x=1cm, y=1cm,
  cause/.style={anchor=east, text width=3.45cm, align=right, inner sep=1pt},
  cat/.style={draw, fill=gray!15, font=\footnotesize\bfseries, inner sep=2pt},
  bone/.style={thick}]
  \newcommand{\cz}[4]{\node[cause,#3] at ({#1-0.2727*abs(#2)-0.12},#2) {#4};}
  % Espina central y efecto
  \draw[very thick,->] (3.3,0) -- (12.4,0);
  \node[draw, very thick, anchor=west, text width=3.2cm, align=center, font=\footnotesize\bfseries, inner sep=3pt] at (12.4,0)
    {24\,h\,11\,min de degradaci\'on y datos inconsistentes};
  % Detonante
  \node[draw=orange!80!black, fill=orange!20, very thick, anchor=east, text width=3.0cm, align=center, font=\footnotesize\bfseries, inner sep=2pt] at (3.2,0)
    {DETONANTE: mantenimiento de \'optica 100G $\to$ corte de 43\,s};
  % Tecnologia (arriba izquierda)
  \draw[bone] (3.5,4.4) -- (4.7,0); \node[cat] at (3.5,4.7) {Tecnolog\'ia};
  \cz{4.7}{1.7}{text=red!70!black,font=\scriptsize\bfseries}{CAUSA RA\'IZ: Orchestrator promueve primarios entre regiones; las apps no toleran esa latencia}
  \cz{4.7}{3.0}{}{Replicaci\'on as\'incrona: escrituras sin replicar en cada sitio}
  \cz{4.7}{3.9}{}{R\'eplicas de lectura con horas de retraso}
  % Procesos (arriba derecha)
  \draw[bone] (8.5,4.4) -- (9.7,0); \node[cat] at (8.5,4.7) {Procesos};
  \cz{9.7}{1.7}{}{Nunca se reconstruy\'o un cl\'uster completo desde respaldo; restauraci\'on remota lenta}
  \cz{9.7}{3.0}{}{Cola de reintentos: unos 200\,000 webhooks expiraron por TTL}
  \cz{9.7}{3.9}{}{Supuestos no probados (sin pruebas de fallo)}
  % Personas y comunicacion (abajo izquierda)
  \draw[bone] (3.5,-4.4) -- (4.7,0); \node[cat] at (3.5,-4.7) {Personas y comunicaci\'on};
  \cz{4.7}{-1.4}{}{Estados solo verde/amarillo/rojo, sin detalle por componente}
  \cz{4.7}{-2.6}{}{Estimaci\'on lineal de recuperaci\'on, err\'onea}
  \cz{4.7}{-3.7}{}{Comunicado tard\'io (blog a las 07:46 UTC)}
  % Entorno (abajo derecha)
  \draw[bone] (8.5,-4.4) -- (9.7,0); \node[cat] at (8.5,-4.7) {Entorno};
  \cz{9.7}{-1.4}{}{Sin arquitectura activo/activo/activo (N+1 por sitio)}
  \cz{9.7}{-2.6}{}{Partici\'on interna de esa magnitud nunca vista}
  \cz{9.7}{-3.7}{}{Pico de tr\'afico al iniciar la jornada}
\end{tikzpicture}}
\caption{Diagrama de Ishikawa del incidente (elaboraci\'on propia a partir de \cite{github2018}).}
\end{figure}
\vspace{-6pt}
\textbf{Prueba de la causa ra\'iz.} Si Orchestrator no hubiera podido promover primarios entre regiones, el corte de 43~s no habr\'ia dejado a las aplicaciones escribiendo a trav\'es del pa\'is durante unos 40 minutos \cite{github2018}. El informe dice que Orchestrator actu\'o ``como estaba configurado''; la mala alineaci\'on fue de configuraci\'on frente a las expectativas de la aplicaci\'on. \todo{Argumentar con sus palabras; contrastar con la clasificaci\'on de causas de \cite{ghosh2022}}.

\section{Impacto cuantificado}
A partir de la cronolog\'ia y del an\'alisis del incidente, el impacto se cuantifica en cinco dimensiones: servicios afectados, trabajo acumulado durante la recuperaci\'on, integridad de los datos, eventos no entregados e impacto econ\'omico y de nivel de servicio. Las cifras provienen del informe post-mortem oficial de GitHub \cite{github2018}; los tiempos transcurridos se calculan desde el inicio del incidente (21 de octubre de 2018, 22:52~UTC).

\begin{itemize}
    \item \textbf{Servicios afectados:} se vieron afectados m\'ultiples cl\'usteres de bases de datos MySQL que almacenan metadatos no Git (como \textit{issues}, \textit{pull requests}, autenticaci\'on y procesamiento en segundo plano), as\'i como la entrega de \textit{webhooks} y la construcci\'on/publicaci\'on de sitios de GitHub Pages. Durante el incidente, los usuarios pudieron observar datos desactualizados o inconsistentes. No obstante, algunas partes de la plataforma permanecieron disponibles; el informe no especifica un n\'umero exacto de cl\'usteres o usuarios afectados \cite{github2018}.

    \item \textbf{Acumulaci\'on de trabajo pendiente:} a las 16:45~UTC del 22 de octubre se registraron m\'as de 5 millones de eventos de \textit{webhooks} y 80\,000 compilaciones de GitHub Pages en cola. La recuperaci\'on requiri\'o restaurar varios terabytes de datos desde copias de seguridad remotas (realizadas cada 4 horas). Adem\'as, las r\'eplicas de lectura presentaron retrasos de varias horas, que se incrementaron con el pico de tr\'afico (13:15~UTC). A las 06:51~UTC se estim\'o una recuperaci\'on en aproximadamente 2 horas, pero la restauraci\'on completa se extendi\'o hasta las 23:03~UTC (16~h 12~min desde esa estimaci\'on); el procesamiento de la cola acumulada tom\'o 6~h 39~min (16:24--23:03~UTC) \cite{github2018}.

    \item \textbf{Integridad y conciliaci\'on de datos:} se produjeron escrituras divergentes entre regiones: unos segundos de escrituras en la costa este no replicadas a la costa oeste y m\'as de 30 minutos de escrituras acumuladas en la costa oeste. En uno de los cl\'usteres m\'as activos se identificaron 954 escrituras no replicadas en la ventana afectada; el informe se\~nala que el total fue relativamente peque\~no. No se report\'o p\'erdida de datos de usuarios, aunque al momento de la publicaci\'on del informe la conciliaci\'on manual segu\'ia en curso \cite{github2018}.

    \item \textbf{Eventos no entregados:} aproximadamente 200\,000 cargas \'utiles (\textit{payloads}) de \textit{webhooks} superaron el tiempo de vida interno (TTL) y fueron descartadas durante el reprocesamiento. Ante esto, GitHub paus\'o temporalmente el procesamiento y aument\'o el TTL. El informe no indica el porcentaje que representan sobre los m\'as de 5 millones de eventos en cola ni afirma que se hayan perdido definitivamente \cite{github2018}.

    \item \textbf{Impacto econ\'omico y nivel de servicio (SLA):} el informe post-mortem oficial no cuantifica p\'erdidas econ\'omicas, compensaciones ni incumplimientos de acuerdos de nivel de servicio (SLA). Por ello, el impacto se presenta en t\'erminos de duraci\'on y alcance: 24~h 11~min de degradaci\'on parcial del servicio, no una ca\'ida total de la plataforma \cite{github2018}.
\end{itemize}

\section{Relaci\'on con la Unidad 1}
\textbf{Transparencia de replicaci\'on y de fallos.} La primera oculta la existencia de varias copias; la segunda, el fallo y la recuperaci\'on de recursos \cite{vansteen2023}. La replicaci\'on qued\'o expuesta cuando los usuarios recibieron datos inconsistentes de r\'eplicas con horas de retraso (22-oct., 11:12 UTC). La de fallos tambi\'en se rompi\'o: los usuarios no pudieron recibir webhooks ni publicar sitios de Pages durante gran parte del incidente \cite{github2018}.

\textbf{CAP: partici\'on real y disponibilidad sacrificada.} El corte de 43~s aisl\'o el centro de datos primario del este de su hub de red; los nodos de Orchestrator del oeste y de la nube formaron qu\'orum y promovieron primarios en el oeste \cite{github2018}. Ante una partici\'on, CAP impide garantizar simult\'aneamente consistencia fuerte (C) y disponibilidad de todas las operaciones (A) \cite{gilbert2002}. Las aplicaciones redirigieron escrituras al oeste después de restablecerse la conexión; esto no demuestra disponibilidad durante el corte. La decisi\'on expl\'icita fue sacrificar disponibilidad para proteger los datos: a las 23:19 UTC se pausaron webhooks y Pages \cite{github2018}. Interpretamos esta medida como una prioridad de consistencia e integridad durante la recuperaci\'on, sin afirmar que todas las lecturas fueran consistentes. La elección C/A durante la partición debe distinguirse de la recuperación posterior \cite{brewer2012}.

\textbf{Tolerancia a fallos y modelo de fallo.} El \emph{failover autom\'atico} con r\'eplicas no mantuvo el servicio esperado: Orchestrator actu\'o seg\'un su configuraci\'on, pero las aplicaciones del este no toleraban la latencia de escribir en el oeste. El consenso de Orchestrator no garantizaba la sincronizaci\'on de las bases MySQL \cite{github2018}. Clasificamos el detonante como \emph{omisi\'on de comunicaci\'on} (interrupci\'on del intercambio entre sitios); sus efectos incluyeron fallos de \emph{temporizaci\'on} (latencia excesiva) y de \emph{respuesta} (valores obsoletos) \cite{vansteen2023}.

\section{Mitigaci\'on aplicada por la empresa}
\textbf{Inmediatas} \cite{github2018}: bloquear el despliegue; estado rojo; pausar webhooks y Pages; restaurar desde respaldos y luego directo desde la costa oeste; aprovisionar r\'eplicas de lectura extra; subir el TTL de los webhooks.\\
\textbf{Permanentes} \cite{github2018} (con la causa del diagrama que atacan): (1) impedir que Orchestrator promueva primarios entre regiones $\to$ causa ra\'iz; (2) nuevo estado por componente $\to$ comunicaci\'on; (3) arquitectura activo/activo/activo $\to$ entorno; (4) probar supuestos y practicar ingenier\'ia del caos $\to$ procesos.

\section{Lecciones aprendidas para el proyecto propio}
\begingroup\scriptsize
\begin{tabularx}{\linewidth}{@{}p{0.4cm}p{1.9cm}XXXX@{}}
\toprule
N. & Componente del PFC & Escenario de fallo & Relaci\'on con el caso & Medida en TA-PFC-E1 & C\'omo se verificar\'a\\
\midrule
1 & Registro de solicitudes del SGRA & Corte entre backend y PostgreSQL; el reintento duplica una solicitud. & Como en GitHub, una falla de red dificulta recuperar escrituras. & Definir transacciones e identificadores únicos para evitar duplicados al reintentar. & Prueba prevista: cortar la conexión y reenviar; debe persistir una sola solicitud.\\
2 & \todo{} & \todo{} & \todo{} & \todo{} & \todo{}\\
3 & \todo{} & \todo{} & \todo{} & \todo{} & \todo{}\\
\bottomrule
\end{tabularx}
\endgroup

\clearpage
% ------------------------- NO CUENTAN PARA LAS 2 PAGINAS -------------------------
\section*{Declaraci\'on de uso de inteligencia artificial}
\small
Se us\'o Claude (Anthropic) para armar la plantilla LaTeX y proponer fuentes.

\printbibliography[title={Referencias}]

\end{document}
